\begin{lemma}[Situation 7.3 is impossible]
Fix a finite loopless multigraph with its individual edge identities, a
palette of \(k\) colors, a distinguished edge \(e_0\), an input coloring,
an exact terminal tree-sequence \(T\), a generating coloring, a
positive-rung extended Tashkinov tree, a hierarchy base, a hierarchy
counterexample \(C\), and a
\noderef{HierarchyMinimumCounterexampleContext} for these data.  Assume
that every vertex has a missing color under every proper coloring leaving
exactly \(e_0\) uncolored.  If
\[
 2\le j:=\noderef{HierarchyPathNumber}(C.\mathit{hierarchy})
   <\ell:=\noderef{HierarchyFullTailLength}(C),
\]
then these data are impossible.
\end{lemma}
\begin{proof}
Let
\[
 I=\noderef{HierarchyEligibleCommonMissingIndices}
\]
for the supplied minimum counterexample.  We first split according as
\(I\) is nonempty or empty.  In the nonempty branch,
\noderef{HierarchySituationThreeClaimJMinusOneNonempty} supplies a
replacement presentation \(P\), a color \(\alpha\), and the full Claim
\(j-1\) package: \(\alpha\) is missing at the terminal vertex and on
\(B=T(y_{j-2})\), and it satisfies the stated outside-or-Gamma
alternative.  In the empty branch,
\noderef{HierarchySituationThreeClaimJMinusOneEmpty} supplies a
replacement presentation \(P\), a color \(\alpha\), and exactly the same
package.  Thus in either branch fix such \(P,\alpha\).

Apply \noderef{HierarchySituationThreeColorChoice} to this full alpha
package.  It supplies a color \(\beta\ne\alpha\), missing on \(B\),
outside the last Gamma union and the exact used-tail color set, and
satisfying precisely the two-alternative color-choice condition needed
below.  Put
\[
 X=T(y_{j-1})=\noderef{HierarchyFullTailPrefix}(C,j-1)
\]
and let
\[
 Q=\noderef{TwoColorComponent}
   (G,P.\mathit{coloring},\alpha,\beta,y_\ell)
\]
be the alpha-beta component rooted at the terminal vertex \(y_\ell\) of
\(T\).

If \(Q\) contains a vertex of \(X\), all hypotheses of
\noderef{HierarchySituationThreeComponentIntersectionShortening} are the
two packages just obtained together with \(2\le j<\ell\) and this
intersection witness.  That lemma gives a contradiction directly.

It remains to consider the case in which \(Q\) contains no vertex of
\(X\).  Negating the existential intersection assertion says explicitly
that
\[
  \operatorname{inl}(v)\notin Q
       \qquad\text{for every }
       v\in\noderef{TreeSequenceVertexSet}(X).
\]
Together with the alpha and beta packages, the path bounds, and the
universal missing-color hypothesis, this is exactly the input to
\noderef{HierarchySituationThreeDisjointPreparation}.  Choose the
resulting datum
\[
 d:\noderef{HierarchySituationThreeDisjointData}(P,\alpha,\beta).
\]
Write \(H=d.\mathit{switched.counterexample.hierarchy}\), let \(q\) be
its last stored index, and put
\[
 D=\noderef{HierarchyStoredColorsAt}(H,q).
\]

If \(d.\theta\in D\), then
\noderef{HierarchySituationThreeDisjointThetaStoredImpossible}, applied
to \(d\) and this membership, gives a contradiction.  Suppose instead
that \(d.\theta\notin D\).  By the definition of
\noderef{HierarchySituationThreeDisjointData}, the field
\(d.\mathit{conditionNine}\) is exactly a disjunction: either the outside
alternative of condition (9) holds, or its stored-beta alternative
holds.  In the first branch,
\noderef{HierarchySituationThreeDisjointThetaUnstoredOutsideImpossible}
applied to \(d\), the theta nonmembership, and the outside package gives
a contradiction.  In the second branch,
\noderef{HierarchySituationThreeDisjointThetaUnstoredBetaStoredImpossible}
applied to \(d\), the same theta nonmembership, and the stored-beta
package gives a contradiction.  These cases exhaust \(I\), the
component intersection decision, theta storedness, and condition (9), so
Situation 7.3 is impossible.
\end{proof}
